% #############################################################################
% Abstract Text
% !TEX root = ../main.tex
% #############################################################################
% reset acronyms
\acresetall
% Guide caps the abstract at 250 words.
% use \noindent in firts paragraph
\noindent Generative Artificial Intelligence and Large Language Models have moved from autocomplete assistants to active contributors across the Software Development Life Cycle, enabling development in which practitioners express requirements in natural language and supervise the resulting artefacts. Adoption, however, has outpaced method. A Systematic Literature Review over 646 articles confirms substantial benefits alongside recurring risks in reliability, security, traceability and accountability, and identifies the absence of structured, repeatable guidance for integrating these tools into professional workflows as the field's central unresolved problem. A majority of functionally correct AI-generated code has been reported to carry security defects, and delivery bottlenecks shift towards review rather than disappearing: value is conditional on process, not on model capability alone.

This dissertation addresses that gap through the Design Science Research Methodology. It proposes a process-oriented framework for governed Human-AI collaboration treating AI systems as first-class collaborators whose contributions must be bounded, validated and traced. The framework defines lifecycle phases with phase-specific quality gates, responsibilities expressed as functions rather than job titles, and governance metrics, anchored by a persistent versioned specification linking every generated artefact to its originating requirement. It is instantiated as Apex, a web application operationalising each phase, gate and artefact against a real project management backend. Evaluation combines application in a real project, standardised usability and workload instruments, interviews with practitioners and Agile experts, and an analytical assessment against criteria set in advance. The contribution is an evidence-grounded process model allowing organisations to adopt AI assistance without surrendering oversight of software quality, security and accountability.
